\odiseachapter{agamenon}{Agamenón y Clitemnestra}\label{cap:agamenon}

Si hubiera que buscar un ejemplo de matrimonio fallido en los griegos, Agamenón y Clitemnestra serían unos candidatos excelentes. Nolan nos da una versión breve, pero muy convincente, de sus problemas conyugales, que resuelven a puñalada limpia. Nada que objetar, por cierto, como tampoco tengo queja del retrato que el biensabido director hace del caudillo aqueo. Incluso el atuendo Darth Vader que le da contribuye a marcar el carácter siniestro del personaje, que se gana el raro honor de ser el tipo más antipático, con diferencia, de la \emph{Ilíada}. En la \emph{Odisea} aparece como una sombra en el Hades ---después de que los dioses, Néstor y Menelao hayan contado su asesinato a manos de Egisto--- y, sobre todo, se dedica a lloriquear.

Vale la pena repetirlo. En la \emph{Odisea}, es Egisto ---no Clitemnestra--- quien lo mata, tendiéndole una emboscada que acaba en escabechina. Ulises lo cuenta así:

\begin{versocita}
de Agamenón el Atrida se me presentaba allí el alma\\
muy apenada; y con ella otras muchas allí se arrimaban\\
de los que en casa de Egisto hallaron con él hora mala.\\
Él enseguida me reconoció, que en la sangre ya andaba,\\
y me lloraba en amargo quejido soltando mil lágrimas,\\
y a mí tendía sus manos como si tocarme anhelara;\\
pero ni aquel firme brío ni aquel su vigor allí estaban,\\
esos que antaño en sus ágiles miembros señoreaban.
\end{versocita}

Nuestro héroe, que no conoce el triste fin del caudillo, se queda de piedra.

\begin{versocita}
Mi corazón se llenó de tristeza y al verlo lloraba,\\
y le decía, alzando la voz de verbos alada:\\
\stanzabreak
«Agamenón, de guerreros señor, gloriosísimo Atrida,\\
¿qué malhadar del morir sombriluengo domó esa tu vida?\\
¿Lo hizo en tus naves quizá Poseidón cuando en ellas volvías\\
y hórrido aliento de pérfido viento te echó encima un día?\\
¿O fueron hombres, quizá, desalmados del mar a la orilla\\
porque sus buenos rebaños de ovejas y bueys les abrías,\\
o porque por su ciudad, o sus hembras quizá, combatías?».
\end{versocita}

En fin, Odiseo enumera las maneras estándar de morir entre los guerreros que regresaban de Troya. O naufragando, o en una refriega, causada por la mala costumbre de los argivos de robar ganado ajeno o asaltar las ciudades que se encontraban de camino, ávidos de botín. La fórmula, de hecho, se repite en la segunda incursión que el lector hace a los infiernos, siguiendo las almas de los pretendientes, en el último canto de la épica. En todo caso, el destino del Atrida ha sido muy distinto al que le plantean. Agamenón contesta:

\begin{versocita}
Sangre Laercia, Odiseo milmañas, hijo del cielo,\\
no me domó Poseidón en las naves a nuestro regreso\\
luego de alzarnos un hórrido aliento de pérfido viento,\\
ni fueron hombres hostiles, ya en tierra, los que me rindieron,\\
que Egisto fue el que dispuso mi sino y momento postrero\\
y me mató con mi odiosa mujer, tras llamarme a su techo,\\
mientras cenaba, tal buey que uno mata en el arrendadero.
\end{versocita}

La escena que relata no puede ser más sangrienta ---a Homero el gore se le da estupendamente---:

\begin{versocita}
Morí la más triste muerte; y, conmigo, mis compañeros,\\
que igual que a cerdos dientirreblancos matanza les dieron\\
como si en casa de un rico señor poderoso con ellos\\
boda o escote o banquete solemne estuvieran haciendo.
\end{versocita}

Se queja Agamenón de que una cosa es caer en combate y otra que te degüellen como una res durante la cena. Tres mil años más tarde, el célebre autor de \emph{Juego de tronos} dejaría al mundo patidifuso con la infame boda roja, una escena que parece salida de versos como estos:

\begin{versocita}
Sé que de muchos guerreros la muerte te vino al encuentro\\
en fragorosa contienda o llevándote a otro en un duelo,\\
mas se te habría deshecho el alma muy más viendo aquello:\\
cómo a redor de las mesas repletas y de los calderos\\
yacíamos y la sangre corría empapando aquel suelo.
\end{versocita}

Tampoco se escapa la pobre Casandra ---que ha profetizado su propia muerte---. Es Clitemnestra quien la mata:

\begin{versocita}
Fue el de la hija de Príamo el frémito más lastimero,\\
el de Casandra, mientras Clitemnestra le hacía degüello\\
sobre mi cuerpo; yo alzaba los brazos, que nada pudieron\\
porque de hierro muriendo yo estaba; y la de ojos de perro\\
no esperó allí ---aunque hacer ya me viera camino al infierno---\\
para sellarme la boca y mis ojos cerrar con sus dedos.
\end{versocita}

Agamenón se refiere a su esposa como \gr{κυνῶπις} (kynôpis), «la de ojos de perro», usando el mismo término que Helena se aplica a sí misma en Esparta. En ambos casos, el juicio moral es menos fuerte de lo que puede parecerle a un lector moderno. El perro homérico es el animal sin vergüenza que te mira directamente a los ojos y, por tanto, la traducción más exacta del epíteto, como ya comentamos, es «descarada». Da que pensar. Parece que a Agamenón le molesta más la desvergüenza de la asesina ---que lo mata sin honor, degollando de paso a Casandra--- que el propio crimen. Los aqueos y su extraño sentido del honor:

\begin{versocita}
No, más horrible no hay nada y tampoco nada más perro\\
que la mujer que acomodo da a un crimen por sus adentros\\
como el que aquella se avino a fraguar, deshonesto y obsceno,\\
de darle muerte a su esposo de ley. Me esperaba yo, ingenuo,\\
ser recibido con alegría por hijos y siervos\\
porque volvía, pero la madama de lutos y duelos\\
sobre ella misma vertió la vergüenza y sobre cuantas luego,\\
pobres mujeres, habrán de venir, aun obrando una recto.
\end{versocita}

Odiseo se compadece, con viril solidaridad, del camarada caído:

\begin{versocita}
Así me habló, y respondiéndole entonces de nuevo le digo:\\
\stanzabreak
«Ay, a la prole de Atreo el buen Zeus clarividino\\
mucho la debe de odiar ---basta ver de esas dos los aliños---\\
desde el principio, que muchos por culpa de Helena caímos,\\
y Clitemnestra te armaba celada, estando tú ido».
\end{versocita}

De ahí que Agamenón le aconseje que de las mujeres no hay que fiarse:

\begin{versocita}
Tal dije, y él enseguida me estaba otra vez respondiendo:\\
\stanzabreak
Por eso nunca ni con tu mujer seas tú un bendituelo\\
de los que de esto y aquello ---de todo--- le dan argumento:\\
cosas le cuentas algunas, mas lleva las más en secreto;
\end{versocita}

También le confirma que, en asunto de mujeres, Odiseo tiene suerte. Aunque lleva un año pasándoselo estupendamente con Circe, Penélope lo espera en casa y allí lo espera también su hijo. Agamenón se lamenta a su vez de no haber podido ver al suyo y es casi el único momento en el que nos compadecemos de él.

\begin{versocita}
y eso que a ti no el morir te vendrá de tu esposa, Odiseo,\\
que es toda juicio y solo consejo de bien lleva dentro\\
la joven hija de Icario, Penélope milpensamientos.\\
Era recién desposada, nosotros ya estábamos yendo\\
para la guerra y allí la dejamos con un niño al pecho,\\
una criatura que ya ocupará entre los hombres asiento,\\
feliz; y, claro, su padre por fin lo verá a su regreso,\\
y correrá él a abrazarse a su padre, como es mandamiento.\\
La mía no permitió, mi mujer, que del hijo que tengo\\
llenos quedaran mis ojos; y muerte me dio antes de eso.
\end{versocita}

Y remata:

\begin{versocita}
Y aún algo más te diré, y procúralo en mientes despierto:\\
secreto, siempre secreto, que no a descubierta, el acerco\\
haz de la nave a tu patria, que yo ya en mujeres no creo.
\end{versocita}

¡Pobre! Qué desengañado está de las féminas. Se le olvidan algunos detallitos, como su afición a apropiarse de toda mujer que se cruza en su camino, empezando por Criseida, la cautiva que se empeña en conservar al principio de la \emph{Ilíada}, a pesar del exuberante rescate que ofrece su padre. No solo se niega a liberar a la muchacha, sino que vocea a quien quiera oírlo que el destino de la pobre chica será envejecer sirviéndolo, literalmente, entre el telar y su cama. De paso, también compara públicamente a su esposa Clitemnestra con Criseida y afirma que es inferior a esta en cuerpo, talle, juicio y labores. Homero no dice que la pobre Casandra ---también cautiva de guerra--- fuera su amante, pero, visto cómo se las gasta el personaje, no cabe ninguna duda. Por otra parte, Egisto lo mata como un animal y Clitemnestra asesina a sangre fría a una pobre inocente. Menuda familia.

Nolan nos cuenta una historia bastante diferente, en la que es Clitemnestra quien mata a su marido ---Casandra ni aparece--- y lo hace por una venganza que, justa o no, cualquier madre, cualquier padre, comprende. Diez años antes, Agamenón ha sacrificado a su hija Ifigenia con el ánimo de que los dioses le concedan buen viento para viajar a Troya\footnote{La historia sigue el \emph{Agamenón} (458 a. C.) de Esquilo y, hasta el desenlace, la \emph{Ifigenia en Áulide} (405 a. C.) de Eurípides.}. Hay, además, un estupendo giro en el film. Lupita Nyong'o interpreta tanto a Helena como a Clitemnestra, algo bastante convincente, ya que las dos reinas son medias hermanas. El espectador se queda, la verdad sea dicha, razonablemente satisfecho con el matarile de Agamenón, y eso que Nolan se salta detalles que en Eurípides ponen los pelos de punta. De hecho, en mi \emph{Abandonando Ítaca} también prefiero la versión de Nolan a la de la \emph{Odisea}. En mi relato, por cierto, el amigo Odiseo no sale muy bien parado. Los remordimientos por esos pecados le consumen en Ítaca.

\begin{versopropio}
La guerra empezó mal,\\
y ahora que eres anciano,\\
doblado como los antiguos olivos de tu isla,\\
te preguntas cómo no lo viste venir,\\
te preguntas cómo fuiste tan estúpido,\\
y tan cruel.
\end{versopropio}

La ausencia de viento hacía imposible partir hacia Troya y los aqueos tenían prisa: había una afrenta que vengar y, sobre todo, una rica ciudad que saquear. Calchas, el adivino, decide que hace falta un sacrificio.

\begin{versopropio}
Ese idiota de Calchas,\\
culpando al capricho de los dioses por la calma chicha,\\
pidiendo una ofrenda de sangre.\\
No una víctima cualquiera,\\
sino que nada menos que Ifigenia, la hija de Agamenón,\\
tenía que ser sacrificada para complacer a las deidades.
\end{versopropio}

El anciano Odiseo, que está de vuelta de todo, sospecha que detrás de la profecía no había dioses, sino una conspiración.

\begin{versopropio}
¿Quién pagaba al sacerdote para que hiciese daño al rey?\\
¿Diomedes, Néstor? Erais muchos, y ninguno de vosotros,\\
---quizás con excepción del obtuso de Áyax y del ingenuo Aquiles---\\
tenía buena fe. ¿Mas matar a una niña?\\
Tal vez Calchas tenía sus propias cuentas pendientes,\\
quizás los dioses estaban realmente aburridos y querían algo de diversión.\\
Salvo que los dioses ---tu vejez lo ha revelado---,\\
solo existen en la imaginación de los hombres,\\
y así, la sangre de la princesa\\
manchó inútilmente vuestras manos.
\end{versopropio}

Pero no todas las manos por igual. Odiseo es casi tan culpable como Agamenón:

\begin{versopropio}
Especialmente tus manos. Fuiste el más insistente,\\
quien declaró delante del resto de los reyes\\
que la vida de una niña, fuese princesa o no,\\
no podía retrasar la empresa. Y fuiste tú\\
quien inventó esa sórdida historia,\\
---que iba a casarse con Aquiles--- para\\
atraer a la pobre víctima a su perdición.
\end{versopropio}

Cierto, los remordimientos le corroen y se lo tiene bien merecido.

\begin{versopropio}
Ahora recuerdas las caras de Ifigenia,\\
la cara que viste cuando irrumpió en la sala,\\
preguntando por su amado padre, creyendo aún\\
que había sido convocada para desposar\\
al mayor guerrero vivo. Oh, ¿recuerdas sus ojos?\\
Brillando como solo brillan las estrellas, sonriéndoos a todos\\
con la inocencia de una muchacha virgen.
\stanzabreak
Y luego, la otra cara, todavía de niña,\\
pero ahora rota y asustada,\\
preguntando a su padre por qué debe morir,\\
y él murmurando mentiras enfermizas,\\
sobre dioses, la gloria y el honor,\\
mientras hunde el cuchillo en su garganta.
\stanzabreak
Una tercera cara te persigue,\\
la de Ifigenia muerta, con sus ojos abiertos en la nada,\\
vacíos de emoción, salvo, quizás,\\
una última sombra de tristeza.\\
Mira la nada y, a un tiempo,\\
te mira fijamente a ti, preguntando: ¿por qué?
\end{versopropio}

¿Por qué la muerte de una niña, de cualquier niña? La Atenea de Nolan murmura algunas frases incoherentes, argumentando que el valor del sacrificio se mide en lo que paga el que lo realiza. ¡Pobre Agamenón, qué disgusto se llevaría pasando a cuchillo a su pobre hija, pobre Odiseo, qué mal le sienta recordar su vileza!

No, el sacrificio de un inocente no admite justificación alguna. A estas alturas, tres mil años después de Homero, deberíamos haber aprendido una lección tan simple que, desafortunadamente ---basta seguir las noticias durante cinco minutos--- se nos sigue escapando.
